\def\cellsz{0.35cm}
\providecommand{\maraGridLineColor}{black!15}
\providecommand{\maraGridBorderColor}{black!30}
\providecommand{\maraGridLineWidth}{0.4pt}
\providecommand{\maraBorderLineWidth}{3pt}
\providecommand{\maraBorderRadius}{2pt}
\providecommand{\TaskInnerTitleFont}{\tiny}
\providecommand{\TaskInnerTitlePlanningX}{-0.6}
\providecommand{\TaskInnerTitleTopY}{6.0}

\begin{scope}[x=\cellsz, y=\cellsz, line cap=round, line join=round]

  \begin{scope}
    % Backdrop panel (reduced top gap, slightly more bottom room)
    \path[fill=black!5] (0,0) rectangle (14.6,10);

    \begin{scope}[local bounding box=allcontent, shift={(1, 2.8)}]
      \begin{scope}[]
        % Background
        \path[fill=\maraGridBorderColor, rounded corners=\maraBorderRadius]
        ([xshift=-\maraBorderLineWidth,yshift=-\maraBorderLineWidth]0,0) rectangle ([xshift=\maraBorderLineWidth,yshift=\maraBorderLineWidth]5,5);

        \fill[black] (0,0) rectangle (5,5);

        % TOP PANEL - Action viewer (mirrors slider example style)
        \begin{scope}
          \def\topPanelY{5.5}
          \def\topPanelHeight{1.2}

          % Top panel background (no border)
          \path[fill=black!20, rounded corners=\maraBorderRadius]
          (0,\topPanelY) rectangle (5,\topPanelY+\topPanelHeight);

          % Action cells
          \begin{scope}[shift={(0.5,\topPanelY+0.2)}]
            \def\actionCellSize{0.8}
            \def\actionSpacing{1.0}

            % Current action (highlighted)
            \fill[gray!60] (0*\actionSpacing,0) rectangle ++(\actionCellSize,\actionCellSize);
            \node[white, font=\tiny] at (0*\actionSpacing+\actionCellSize/2,\actionCellSize/2) {$\bullet$};

            % Future actions (dimmed)
            \fill[black!60] (1*\actionSpacing,0) rectangle ++(\actionCellSize,\actionCellSize);
            \node[white, font=\tiny] at (1*\actionSpacing+\actionCellSize/2,\actionCellSize/2) {---};

            \fill[black!60] (2*\actionSpacing,0) rectangle ++(\actionCellSize,\actionCellSize);
            \node[white, font=\tiny] at (2*\actionSpacing+\actionCellSize/2,\actionCellSize/2) {---};

            \fill[black!60] (3*\actionSpacing,0) rectangle ++(\actionCellSize,\actionCellSize);
            \node[white, font=\tiny] at (3*\actionSpacing+\actionCellSize/2,\actionCellSize/2) {---};

            % Grid lines for action cells
            \draw[draw=\maraGridLineColor, line width=\maraGridLineWidth]
            (0,0) -- (4*\actionSpacing,0) -- (4*\actionSpacing,\actionCellSize) -- (0,\actionCellSize) -- cycle;
            \foreach \i in {1,2,3} {
                \draw[draw=\maraGridLineColor, line width=\maraGridLineWidth]
                (\i*\actionSpacing,0) -- (\i*\actionSpacing,\actionCellSize);
              }
          \end{scope}
        \end{scope}

        % Colored cells
        \begin{scope}[fill={rgb,1:red,0.000;green,0.000;blue,1.000}, fill opacity=1.000]
          \fill (0,4) rectangle ++(1,1);
        \end{scope}
        % \begin{scope}[fill={rgb,1:red,0.000;green,0.000;blue,0.000}, fill opacity=1.000]
        %   \fill (1,4) rectangle ++(1,1);
        %   \fill (2,4) rectangle ++(1,1);
        %   \fill (3,4) rectangle ++(1,1);
        %   \fill (4,4) rectangle ++(1,1);
        %   \fill (1,3) rectangle ++(1,1);
        %   \fill (4,0) rectangle ++(1,1);
        % \end{scope}
        \begin{scope}[fill={rgb,1:red,0.502;green,0.502;blue,0.502}, fill opacity=1.000]
          \fill (0,3) rectangle ++(1,1);
        \end{scope}
        \begin{scope}[fill={rgb,1:red,1.000;green,0.843;blue,0.000}, fill opacity=1.000]
          \fill (1,2) rectangle ++(1,1);
        \end{scope}


        % (no crosshatch on left grid after swapping)




        % Grid lines - Vertical
        \begin{scope}[draw=\maraGridLineColor, line width=\maraGridLineWidth, line cap=round]
          \draw (0,0) -- (0,5);
          \draw (1,0) -- (1,5);
          \draw (2,0) -- (2,5);
          \draw (3,0) -- (3,5);
          \draw (4,0) -- (4,5);
          \draw (5,0) -- (5,5);
        \end{scope}

        % Grid lines - Horizontal
        \begin{scope}[draw=\maraGridLineColor, line width=\maraGridLineWidth, line cap=round]
          \draw (0,0) -- (5,0);
          \draw (0,1) -- (5,1);
          \draw (0,2) -- (5,2);
          \draw (0,3) -- (5,3);
          \draw (0,4) -- (5,4);
          \draw (0,5) -- (5,5);
        \end{scope}

        % Highlight box on left grid (cover entire unshaded region)
        \node[inner sep=0pt, outer sep=0pt, draw=none, fit={(0,0) (2,5)}] (PLHighlightBox) {};
        \draw[blue!60, line width=1.2pt, rounded corners=1pt]
        (PLHighlightBox.south west) rectangle (PLHighlightBox.north east);

        \node[
          draw=\maraGridBorderColor,
          fill=black!40,
          rounded corners=\maraBorderRadius,
          inner sep=1.5pt,
          minimum width=1.6cm,
          minimum height=0.35cm,
          font=\tiny,
          text=black
        ] at (2.5,-1.5) {Give up};
      \end{scope}

      % Second (right) grid (slightly smaller)
      \begin{scope}[shift={(7,0)}]
        % Label above the right grid (two lines, left-aligned with grid)
        \node[font=\TaskInnerTitleFont, text=black, align=left, anchor=west] at (\TaskInnerTitlePlanningX, \TaskInnerTitleTopY) {Reach the target\\configuration:};
        % Outer frame
        \path[fill=\maraGridBorderColor, rounded corners=\maraBorderRadius]
        ([xshift=-\maraBorderLineWidth,yshift=-\maraBorderLineWidth]0,0)
        rectangle ([xshift=\maraBorderLineWidth,yshift=\maraBorderLineWidth]5,5);
        % Background
        \fill[black] (0,0) rectangle (5,5);

        % Colored cells (same layout)
        \begin{scope}[fill={rgb,1:red,0.000;green,0.000;blue,1.000}, fill opacity=1.000]
          \fill (0,4) rectangle ++(1,1);
        \end{scope}
        % \begin{scope}[fill={rgb,1:red,0.000;green,0.000;blue,0.000}, fill opacity=1.000]
        %   \fill (1,4) rectangle ++(1,1);
        %   \fill (2,4) rectangle ++(1,1);
        %   \fill (3,4) rectangle ++(1,1);
        %   \fill (4,4) rectangle ++(1,1);
        %   \fill (1,3) rectangle ++(1,1);
        %   \fill (4,0) rectangle ++(1,1);
        % \end{scope}
        \begin{scope}[fill={rgb,1:red,0.502;green,0.502;blue,0.502}, fill opacity=1.000]
          \fill (0,1) rectangle ++(1,1);
        \end{scope}
        \begin{scope}[fill={rgb,1:red,1.000;green,0.843;blue,0.000}, fill opacity=1.000]
          \fill (1,2) rectangle ++(1,1);
        \end{scope}

        % Grid lines - Vertical
        \begin{scope}[draw=\maraGridLineColor, line width=\maraGridLineWidth, line cap=round]
          \draw (0,0) -- (0,5);
          \draw (1,0) -- (1,5);
          \draw (2,0) -- (2,5);
          \draw (3,0) -- (3,5);
          \draw (4,0) -- (4,5);
          \draw (5,0) -- (5,5);
        \end{scope}

        % Grid lines - Horizontal
        \begin{scope}[draw=\maraGridLineColor, line width=\maraGridLineWidth, line cap=round]
          \draw (0,0) -- (5,0);
          \draw (0,1) -- (5,1);
          \draw (0,2) -- (5,2);
          \draw (0,3) -- (5,3);
          \draw (0,4) -- (5,4);
          \draw (0,5) -- (5,5);
        \end{scope}

        % White highlight box on right grid
        % Apply crosshatch (masked region) on right grid
        % Function to draw crosshatch patterns
        \newcommand{\drawCrosshatch}[2]{
          \begin{scope}[shift={(#1,#2)}]
            % Main diagonal lines
            \draw[gray, line width=0.8pt]
            (0.1,0.1) -- (0.9,0.9);
            \draw[gray, line width=0.8pt]
            (0.1,0.9) -- (0.9,0.1);
            % Additional crosshatch lines for density
            \draw[gray, line width=0.8pt]
            (0.1,0.3) -- (0.7,0.9);
            \draw[gray, line width=0.8pt]
            (0.3,0.1) -- (0.9,0.7);
            \draw[gray, line width=0.8pt]
            (0.1,0.7) -- (0.3,0.9);
            \draw[gray, line width=0.8pt]
            (0.7,0.1) -- (0.9,0.3);
            \draw[gray, line width=0.8pt]
            (0.1,0.5) -- (0.5,0.9);
            \draw[gray, line width=0.8pt]
            (0.5,0.1) -- (0.9,0.5);
          \end{scope}
        }

        \drawCrosshatch{4}{0}
        \drawCrosshatch{4}{1}
        \drawCrosshatch{4}{2}
        \drawCrosshatch{4}{3}
        \drawCrosshatch{4}{4}
        \drawCrosshatch{3}{0}
        \drawCrosshatch{3}{1}
        \drawCrosshatch{3}{2}
        \drawCrosshatch{3}{3}
        \drawCrosshatch{3}{4}
        \drawCrosshatch{2}{0}
        \drawCrosshatch{2}{1}
        \drawCrosshatch{2}{2}
        \drawCrosshatch{2}{3}
        \drawCrosshatch{2}{4}

        % Highlight box on right grid (cover entire unshaded region)
        \draw[blue!60, line width=1.2pt, rounded corners=1pt]
        (0,0) rectangle (2,5);
      \end{scope}
    \end{scope}


  \end{scope}
\end{scope}
